% GENERATED FILE. Edit canonical lessons, then run npm run generate:books.
% canonical-source-hash: fnv1a64:33c4dad5bd603c33
% canonical-lessons: MR-C341-kadambari, MR-C341-caritra, MR-C341-adbhut, MR-C341-avismaraniy, MR-C341-ciradne, MR-R341-review-335-341, MR-R341-review-final-341

\chapter{Books to Remember}
\label{ch:mr-books-to-remember}
\hlchaptermodality{341}

\begin{chapteropening}
\textbf{By the end of this chapter:} \emph{I can talk about novels, biographies and unforgettable books.}

Complete the last lesson of chapter 341: I can talk about novels, biographies and unforgettable books.
\end{chapteropening}

\section[\texorpdfstring{\mr{कादंबरी} (kādambarī)}{कादंबरी (kādambarī)}]{\mr{कादंबरी} (kādambarī) — a novel}

\label{lesson:MR-C341-kadambari}



\begin{quote}
Before the new one: say the Marathi for thrilling, then the Marathi for to wrestle with.
\end{quote}

\subsection*{You'll want to know: \mr{कादंबरी}}
\textbf{\mr{कादंबरी}} — \emph{kādambarī} — \textquotedblleft{}a novel\textquotedblright{}.

\begin{quote}
\mr{मी} \mr{सुट्टीत} \mr{एक} \mr{मोठी} \mr{कादंबरी} \mr{वाचली}.
\end{quote}

\begin{quote}
\emph{mī suṭṭīt ek moṭhī kādambarī vāclī.}
\end{quote}

\begin{quote}
I read a long novel during the holidays.
\end{quote}

Feminine.

\mr{कादंबरीकार}, a novelist.

\begin{cousinweb}
\textbf{\mr{कादंबरी}} (\emph{kādambarī}) comes from Sanskrit \emph{Kādambarī}, the title of a celebrated prose romance (Monier-Williams); Marathi made the title the word for every novel.
\end{cousinweb}

\begin{grammarlens}[title={What you've built}]
One more word for telling what happened.
\end{grammarlens}

\subsection*{Your turn}
\begin{itemize}
\raggedright
  \item \emph{Say it:} \emph{kādambarī}
  \item \emph{Say it:} \emph{kādambarī}, once more
  \item \emph{Say it:} \emph{mī suṭṭīt ek moṭhī kādambarī vāclī.} — I read a long novel during the holidays.
  \item \emph{Recall:} say the Marathi for thrilling, then the Marathi for to wrestle with, then say \emph{kādambarī} again
\end{itemize}

\begin{tcolorbox}[breakable,colback=teal!4,colframe=teal!35!black,title={Before you move on}]
What does \mr{कादंबरी} mean? (\textquotedblleft{}A novel\textquotedblright{}.) Say it once more.
\end{tcolorbox}
\section[\texorpdfstring{\mr{चरित्र} (caritra)}{चरित्र (caritra)}]{\mr{चरित्र} (caritra) — a biography}

\label{lesson:MR-C341-caritra}



\begin{quote}
Before the new one: say the Marathi for to wrestle with, then the Marathi for a novel.
\end{quote}

\subsection*{You'll want to know: \mr{चरित्र}}
\textbf{\mr{चरित्र}} — \emph{caritra} — \textquotedblleft{}a biography\textquotedblright{}.

\begin{quote}
\mr{तिने} \mr{एका} \mr{खेळाडूचं} \mr{चरित्र} \mr{लिहिलं}.
\end{quote}

\begin{quote}
\emph{tine ekā kheḷāḍūcaṁ caritra lihilaṁ.}
\end{quote}

\begin{quote}
She wrote the biography of a sportsman.
\end{quote}

Neuter.

\mr{चरित्र} \mr{लिहिणे}, to write a life story.

\begin{cousinweb}
\textbf{\mr{चरित्र}} (\emph{caritra}) is a tatsama, taken from Sanskrit \emph{caritra}, 'behaviour, deeds, adventures' (Monier-Williams).
\end{cousinweb}

\begin{grammarlens}[title={What you've built}]
One more word for telling what happened.
\end{grammarlens}

\subsection*{Your turn}
\begin{itemize}
\raggedright
  \item \emph{Say it:} \emph{caritra}
  \item \emph{Say it:} \emph{caritra}, once more
  \item \emph{Say it:} \emph{tine ekā kheḷāḍūcaṁ caritra lihilaṁ.} — She wrote the biography of a sportsman.
  \item \emph{Recall:} say the Marathi for to wrestle with, then the Marathi for a novel, then say \emph{caritra} again
\end{itemize}

\begin{tcolorbox}[breakable,colback=teal!4,colframe=teal!35!black,title={Before you move on}]
What does \mr{चरित्र} mean? (\textquotedblleft{}A biography\textquotedblright{}.) Say it once more.
\end{tcolorbox}
\section[\texorpdfstring{\mr{अद्भुत} (adbhut)}{अद्भुत (adbhut)}]{\mr{अद्भुत} (adbhut) — wondrous; marvellous}

\label{lesson:MR-C341-adbhut}



\begin{quote}
Before the new one: say the Marathi for a novel, then the Marathi for a biography.
\end{quote}

\subsection*{You'll want to know: \mr{अद्भुत}}
\textbf{\mr{अद्भुत}} — \emph{adbhut} — \textquotedblleft{}wondrous; marvellous\textquotedblright{}.

\begin{quote}
\mr{पुस्तकातलं} \mr{जग} \mr{अद्भुत} \mr{होतं}.
\end{quote}

\begin{quote}
\emph{pustakātlaṁ jag adbhut hotaṁ.}
\end{quote}

\begin{quote}
The world in the book was wondrous.
\end{quote}

\mr{अद्भुत} keeps one form: \mr{अद्भुत} \mr{कथा}, \mr{अद्भुत} \mr{दृश्य}.

\begin{cousinweb}
\textbf{\mr{अद्भुत}} (\emph{adbhut}) is a tatsama, taken from Sanskrit \emph{adbhuta}, 'wonderful, marvellous' (Monier-Williams).
\end{cousinweb}

\begin{grammarlens}[title={What you've built}]
One more word for telling what happened.
\end{grammarlens}

\subsection*{Your turn}
\begin{itemize}
\raggedright
  \item \emph{Say it:} \emph{adbhut}
  \item \emph{Say it:} \emph{adbhut}, once more
  \item \emph{Say it:} \emph{pustakātlaṁ jag adbhut hotaṁ.} — The world in the book was wondrous.
  \item \emph{Recall:} say the Marathi for a novel, then the Marathi for a biography, then say \emph{adbhut} again
\end{itemize}

\begin{tcolorbox}[breakable,colback=teal!4,colframe=teal!35!black,title={Before you move on}]
What does \mr{अद्भुत} mean? (\textquotedblleft{}Wondrous; marvellous\textquotedblright{}.) Say it once more.
\end{tcolorbox}
\section[\texorpdfstring{\mr{अविस्मरणीय} (avismaraṇīy)}{अविस्मरणीय (avismaraṇīy)}]{\mr{अविस्मरणीय} (avismaraṇīy) — unforgettable}

\label{lesson:MR-C341-avismaraniy}



\begin{quote}
Before the new one: say the Marathi for a biography, then the Marathi for wondrous.
\end{quote}

\subsection*{You'll want to know: \mr{अविस्मरणीय}}
\textbf{\mr{अविस्मरणीय}} — \emph{avismaraṇīy} — \textquotedblleft{}unforgettable\textquotedblright{}.

\begin{quote}
\mr{ती} \mr{सहल} \mr{माझ्यासाठी} \mr{अविस्मरणीय} \mr{होती}.
\end{quote}

\begin{quote}
\emph{tī sahal mājhyāsāṭhī avismaraṇīy hotī.}
\end{quote}

\begin{quote}
That trip was unforgettable for me.
\end{quote}

\mr{अविस्मरणीय} keeps one form; \mr{अ}-, not, + \mr{विस्मरणीय}, to be forgotten.

\begin{cousinweb}
\textbf{\mr{अविस्मरणीय}} (\emph{avismaraṇīy}) is built from the Sanskrit prefix \emph{a-}, 'not', and \emph{vismaraṇīya}, 'to be forgotten' (Monier-Williams).
\end{cousinweb}

\begin{grammarlens}[title={What you've built}]
One more word for telling what happened.
\end{grammarlens}

\subsection*{Your turn}
\begin{itemize}
\raggedright
  \item \emph{Say it:} \emph{avismaraṇīy}
  \item \emph{Say it:} \emph{avismaraṇīy}, once more
  \item \emph{Say it:} \emph{tī sahal mājhyāsāṭhī avismaraṇīy hotī.} — That trip was unforgettable for me.
  \item \emph{Recall:} say the Marathi for a biography, then the Marathi for wondrous, then say \emph{avismaraṇīy} again
\end{itemize}

\begin{tcolorbox}[breakable,colback=teal!4,colframe=teal!35!black,title={Before you move on}]
What does \mr{अविस्मरणीय} mean? (\textquotedblleft{}Unforgettable\textquotedblright{}.) Say it once more.
\end{tcolorbox}
\section[\texorpdfstring{\mr{चिरडणे} (ciraḍṇe)}{चिरडणे (ciraḍṇe)}]{\mr{चिरडणे} (ciraḍṇe) — to crush}

\label{lesson:MR-C341-ciradne}



\begin{quote}
Before the new one: say the Marathi for wondrous, then the Marathi for unforgettable.
\end{quote}

\subsection*{You'll want to know: \mr{चिरडणे}}
\textbf{\mr{चिरडणे}} — \emph{ciraḍṇe} — \textquotedblleft{}to crush\textquotedblright{}.

\begin{quote}
\mr{गोष्टीत} \mr{हत्तीने} \mr{राक्षसाला} \mr{पायाखाली} \mr{चिरडलं}.
\end{quote}

\begin{quote}
\emph{goṣṭīt hattīne rākṣasālā pāyākhālī ciraḍlaṁ.}
\end{quote}

\begin{quote}
In the story the elephant crushed the ogre underfoot.
\end{quote}

Past, with -\mr{ने} (he …): \mr{त्याने} \mr{चिरडलं}, \emph{tyāne ciraḍlaṁ}.

Transitive: the one crushed takes -\mr{ला}.

\begin{cousinweb}
\textbf{\mr{चिरडणे}} (\emph{ciraḍṇe}) is a Marathi verb of crushing; its origin is not settled.
\end{cousinweb}

\begin{grammarlens}[title={What you've built}]
One more word for telling what happened.
\end{grammarlens}

\subsection*{Your turn}
\begin{itemize}
\raggedright
  \item \emph{Say it:} \emph{ciraḍṇe}
  \item \emph{Say it:} \emph{ciraḍṇe}, once more
  \item \emph{Say it:} \emph{goṣṭīt hattīne rākṣasālā pāyākhālī ciraḍlaṁ.} — In the story the elephant crushed the ogre underfoot.
  \item \emph{Recall:} say the Marathi for wondrous, then the Marathi for unforgettable, then say \emph{ciraḍṇe} again
\end{itemize}

\begin{tcolorbox}[breakable,colback=teal!4,colframe=teal!35!black,title={Before you move on}]
What does \mr{चिरडणे} mean? (\textquotedblleft{}To crush\textquotedblright{}.) Say it once more.
\end{tcolorbox}
\section[(dialogue)]{Looking back: legends, ballads and books}

\label{lesson:MR-R341-review-335-341}



\begin{quote}
Say the Marathi for an excavation and remains.
\end{quote}

\subsection*{Your turn}
Say these aloud:

\begin{itemize}
\raggedright
  \item an excavation, remains, age-old, to dig, to carve
  \item rock-cut caves, a cave, secret, an underground passage, a treasure
  \item a legend, a heroic ballad, immortal, renowned, a poet
  \item an ogre, a fairy, magic, a witch, a curse
  \item a protagonist, a villain, a character, a holy man, a moneylender
  \item a mystery, the moral, extraordinary, thrilling, to wrestle with
  \item a novel, a biography, wondrous, unforgettable, to crush
\end{itemize}

\begin{tcolorbox}[breakable,colback=teal!4,colframe=teal!35!black,title={Before you move on}]
An excavation? (\textbf{\mr{उत्खनन}}, \emph{utkhanan}.) Remains? (\textbf{\mr{अवशेष}}, \emph{avaśeṣ}.) Age-old? (\textbf{\mr{पुरातन}}, \emph{purātan}.) To dig? (\textbf{\mr{खोदणे}}, \emph{khodṇe}.) To carve? (\textbf{\mr{कोरणे}}, \emph{korṇe}.) Rock-cut caves? (\textbf{\mr{लेणी}}, \emph{leṇī}.) A cave? (\textbf{\mr{गुहा}}, \emph{guhā}.) Secret? (\textbf{\mr{गुप्त}}, \emph{gupta}.) An underground passage? (\textbf{\mr{भुयार}}, \emph{bhuyār}.) A treasure? (\textbf{\mr{खजिना}}, \emph{khajinā}.) A legend? (\textbf{\mr{दंतकथा}}, \emph{dantkathā}.) A heroic ballad? (\textbf{\mr{पोवाडा}}, \emph{povāḍā}.) Immortal? (\textbf{\mr{अमर}}, \emph{amar}.) Renowned? (\textbf{\mr{प्रख्यात}}, \emph{prakhyāt}.) A poet? (\textbf{\mr{कवी}}, \emph{kavī}.) An ogre? (\textbf{\mr{राक्षस}}, \emph{rākṣas}.) A fairy? (\textbf{\mr{परी}}, \emph{parī}.) Magic? (\textbf{\mr{जादू}}, \emph{jādū}.) A witch? (\textbf{\mr{चेटकीण}}, \emph{ceṭkīṇ}.) A curse? (\textbf{\mr{शाप}}, \emph{śāp}.) A protagonist? (\textbf{\mr{नायक}}, \emph{nāyak}.) A villain? (\textbf{\mr{खलनायक}}, \emph{khalnāyak}.) A character? (\textbf{\mr{पात्र}}, \emph{pātra}.) A holy man? (\textbf{\mr{साधू}}, \emph{sādhū}.) A moneylender? (\textbf{\mr{सावकार}}, \emph{sāvkār}.) A mystery? (\textbf{\mr{रहस्य}}, \emph{rahasya}.) The moral? (\textbf{\mr{तात्पर्य}}, \emph{tātparya}.) Extraordinary? (\textbf{\mr{विलक्षण}}, \emph{vilakṣaṇ}.) Thrilling? (\textbf{\mr{थरारक}}, \emph{tharārak}.) To wrestle with? (\textbf{\mr{झगडणे}}, \emph{jhagaḍṇe}.) A novel? (\textbf{\mr{कादंबरी}}, \emph{kādambarī}.) A biography? (\textbf{\mr{चरित्र}}, \emph{caritra}.) Wondrous? (\textbf{\mr{अद्भुत}}, \emph{adbhut}.) Unforgettable? (\textbf{\mr{अविस्मरणीय}}, \emph{avismaraṇīy}.) To crush? (\textbf{\mr{चिरडणे}}, \emph{ciraḍṇe}.)
\end{tcolorbox}
\section[(dialogue)]{Books to Remember — once more}

\label{lesson:MR-R341-review-final-341}



\begin{quote}
Say the Marathi for a novel and a biography.
\end{quote}

\subsection*{Your turn}
\begin{itemize}
\raggedright
  \item \emph{Say it:} a novel, a biography, wondrous, unforgettable, to crush
\end{itemize}

\begin{tcolorbox}[breakable,colback=teal!4,colframe=teal!35!black,title={Before you move on}]
A novel? (\textbf{\mr{कादंबरी}}, \emph{kādambarī}.) A biography? (\textbf{\mr{चरित्र}}, \emph{caritra}.) Wondrous? (\textbf{\mr{अद्भुत}}, \emph{adbhut}.) Unforgettable? (\textbf{\mr{अविस्मरणीय}}, \emph{avismaraṇīy}.) To crush? (\textbf{\mr{चिरडणे}}, \emph{ciraḍṇe}.)
\end{tcolorbox}
